\section{Fourteen positive unknowns at degree thirty}
\label{sec:sc-fourteen}

The construction in this section combines a free Pell ordinate with a new
adaptive quotient modulus.  It gives a fixed polynomial with fourteen
strictly positive integer unknowns, exact total degree $30$, and a
division-free evaluator using $83$ binary arithmetic operations.  The
principal improvement is the modulus $U=Y^2$: it has polynomial degree two,
but a preceding integrality argument forces it to be large enough for the
index recovery.  The proof establishes that size before invoking the
strong Pell-rank theorem.

\begin{theorem}\label{thm:sc-main}
Let $P$ be the polynomial defined in
\eqref{eq:sc-registers}--\eqref{eq:sc-polynomial}, with the complete fixed
coefficient tuple specified below.  Then
\[
 \mathrm{ZFC}\vdash
 \left(\exists\boldsymbol v\in\Z_{>0}^{14}\ P(\boldsymbol v)=0\right)
 \ \longleftrightarrow\
 \left(\text{there is a finite syntactic ZFC refutation}\right).
\]
Consequently
\[
 \mathrm{ZFC}\vdash
 \left(\forall\boldsymbol v\in\Z_{>0}^{14}\ P(\boldsymbol v)\ne0\right)
 \ \longleftrightarrow\ \Con(\mathrm{ZFC}).
\]
The exact degree is $30$.  The supplied arithmetic circuit has $44$
multiplications, $23$ additions, and $16$ subtractions.
\end{theorem}

Here and below, the semantic conclusion composes the new elementary
reduction with the inherited native compiler and the literal
K215/refutation calibration.  The full ordinary mathematical argument and
its source dependencies are supplied; an additional end-to-end proof
assistant certificate is not claimed.  No consistency assumption enters
either direction.

All artifact paths in this section are relative to \path{positive14/}.

\subsection{The fixed polynomial and its coefficient binding}

The fourteen supplied coordinates are
\[
 (J,F,\alpha,z,f,i,T,Y,w,\rho,\eta,\zeta,e,\gamma)
       \in\Z_{>0}^{14}.
\]
All further letters in the definition are polynomial abbreviations, not
additional quantified variables.  Use the \emph{entire} compact native
coefficient tuple from package 44.  Its four ports are
\[
\begin{gathered}
 B=2^{590069},\qquad
 \beta=B-1,\qquad K_0=D_C+BD_R,\\
 M_C=M_{C,\mathrm{native}},\qquad
 M_F=M_{F,\mathrm{native}}+B-1.
\end{gathered}
\]
In particular, $M_F$ is the already shifted source port.  The exact
integers are the unsigned little-endian byte streams obtained by
decompressing the following four delivered files.  Their hashes, finite
coefficient recipes, and binding are recorded in
\path{native/COEFFICIENT_BINDING.json}.
\begin{center}
\begin{tabular}{lll}
\hline
Port & Exact payload under \path{native/coefficients/} & Raw bytes\\
\hline
$\beta$ & \path{beta.bin.gz} & $73\,759$\\
$K_0$ & \path{K.bin.gz} & $74\,111$\\
$M_C$ & \path{MC.bin.gz} & $73\,759$\\
$M_F$ & \path{MF_source.bin.gz} & $73\,759$\\
\hline
\end{tabular}
\end{center}
The coefficient recipe uses radix $2^{67}$ and length $8807$, with the
ordered literal list of $943$ Wang tiles and the $430$ literal K215
transition rows included in \path{native/}.  The independent generator and
auditor are included there too.  Thus this is one fixed four-tuple, rather
than four freely adjustable parameters.  In particular it satisfies
\[
 0<M_C<B-1,\qquad B-1<M_F<2(B-1),
 \quad M_C\equiv2\pmod4,
 \quad M_{F,\mathrm{native}}\equiv4\pmod8,
\]
and the native mask populations sum to $590069$.  Its complete native
interface, including canonical completeness, is proved in
\path{native/PROOF.md}.  The new construction binds all four ports
together to this interface.

Define
\begin{equation}\label{eq:sc-registers}
\begin{aligned}
 q&=\beta J+1,                 & C&=q-F-\alpha,\\
 X&=wq,                       & E&=XY,\\
 k&=\eta+\zeta,               & c&=kY+\eta,\\
 a&=E+Y,                      & H&=4a+3,\\
 \Delta&=a^2+H=(a+2)^2-1,     & \ell&=EkY,\\
 b_0&=X+ac,                   & G&=1+\Delta c^2,\\
 R&=[(q-1)(q-F)+\alpha](q^2-1)+(M_C+qM_F)J,\\
 r&=k-R-1,                    & U&=Y^2,\\
 D_*&=Ub_0+H(U\gamma-r),       & M&=D_*^2-U^2G,\\
 I&=ic^2,                     & A_{\mathrm{aux}}&=\Delta(f^2-1),\\
 W&=I^2-A_{\mathrm{aux}},      & V&=f^2(cT-R)-c,\\
 y&=V+e,                      & N_a&=A_{\mathrm{aux}}(V^2-y^2)+y^2,\\
 L_t&=(K_0+w)C-F-z(q-1),\\
 A_0&=[q^3(\ell+\rho)+Y]^2-q^6[\ell(\ell+k)+1].
\end{aligned}
\end{equation}
The equation is
\begin{equation}\label{eq:sc-polynomial}
 \boxed{P=A_0^2+M^2+W^2+N_a^2+L_t^2-1=0.}
\end{equation}
Only fixed nonnegative powers occur.  The definition contains no division,
variable exponent, recurrence, bit test, or radical.  The literal evaluator
is \path{polynomial.json}; \path{build.py} reconstructs it.  Positivity is
part of the quantified domain in Theorem~\ref{thm:sc-main}.

\subsection{The inherited interface used in the proof}

For an integer $A\ge2$, define the Pell polynomials by
\[
 \chi_A(0)=1,\quad\psi_A(0)=0,\qquad
 \begin{aligned}
 \chi_A(j+1)&=A\chi_A(j)+(A^2-1)\psi_A(j),\\
 \psi_A(j+1)&=\chi_A(j)+A\psi_A(j).
 \end{aligned}
\]
These sequences appear only in the proof.  We use the standard Pell
classification, their strong divisibility, and the elementary recurrence
congruences.  The source proofs are preserved in
\path{upstream/PELL_LEMMAS.md} and
\path{upstream/ORDINARY_KERNEL.md}.

The particular compiler interface needed here is as follows.  Take fifteen
positive coordinates
\[
 (J,F,\alpha,z_{\rm old},f,h,i,T_{\rm old},s,w,
       \tau,\eta,\zeta,y_{\rm old},\gamma_{\rm old}),
\]
put $Y=sq^3$, use the shared abbreviations above, and set
\[
 D=b_0+H\gamma_{\rm old},\qquad
 V_{\rm old}=c(T_{\rm old}f-1)-Rf^2.
\]
The ordinary strong kernel consists of the six equations
\begin{equation}\label{eq:sc-source}
\begin{aligned}
 \tau^2-\ell(\ell+k)&=1, & D^2-\Delta c^2&=1,\\
 A_{\rm aux}(V_{\rm old}^2-y_{\rm old}^2)+y_{\rm old}^2&=1,
       & k-hE-R&=1,\\
 (K_0+w)C+(q-F)-z_{\rm old}(q-1)&=1,
       & (ic^2)^2&=\Delta(f^2-1).
\end{aligned}
\end{equation}
For an authentic complete tuple, a positive solution of this system gives
a finite blank K215 computation, hence a finite ZFC refutation.  Conversely
a finite refutation supplies canonical normalized compiler witnesses, which
give ordinary witnesses by
\begin{equation}\label{eq:sc-normalized}
 i=\Delta i_{\rm norm},\qquad
 f^2=1+\Delta i_{\rm norm}^2c^4.
\end{equation}
The canonical outer identities are
\begin{equation}\label{eq:sc-canonical}
 p=R=2n-1,\qquad R\equiv3\pmod4,\qquad
 X=wq=2^R,\quad Y=sq^3,\quad c=\psi_{a+2}(R).
\end{equation}
These statements are the ordinary compiler theorem in
\path{upstream/ORDINARY_KERNEL.md}, composed with the complete native
adapter in \path{native/PROOF.md}.  They specify the imported semantic
dependency precisely.  Sections below prove the new reduction to and from
this interface.

\subsection{A minimum block replaces an extra squared residual}

\begin{lemma}\label{lem:sc-minimum}
For every positive assignment to the fourteen coordinates,
\[
 P=0\quad\Longleftrightarrow\quad
 A_0=M=W=L_t=0\quad\hbox{and}\quad N_a=1.
\]
\end{lemma}
\begin{proof}
We have $a>0$, $\Delta>0$, $c\ge3$, and $I=ic^2\ge9$.
If $f=1$, then $A_{\rm aux}=0$ and $W=I^2>1$, which is incompatible
with $P=0$.  If $f\ge2$, then $A_{\rm aux}>1$.  In this case $N_a=0$
would imply
\[
 A_{\rm aux}V^2=(A_{\rm aux}-1)(V+e)^2.
\]
Neither $V$ nor $V+e$ could vanish, since simultaneous vanishing would
force $e=0$.  Thus $A_{\rm aux}(A_{\rm aux}-1)$ would be a rational
square.  A rational square root of an integer is integral, whereas
\[
 (A_{\rm aux}-1)^2
 <A_{\rm aux}(A_{\rm aux}-1)<A_{\rm aux}^2.
\]
This contradiction shows $N_a^2\ge1$.  The sum in
\eqref{eq:sc-polynomial} therefore forces the other four residuals to
vanish and $N_a=\pm1$.

With $W=0$ we have $A_{\rm aux}=I^2$, and
\[
 N_a=I^2V^2-(I^2-1)y^2.
\]
Modulo four this is $V^2$ if $I$ is odd and $y^2$ if $I$ is even.
It cannot be $-1$.  The converse follows by substitution.
\end{proof}

\subsection{Scale recovery and the adaptive modulus}

Suppose $P=0$.  From $A_0=0$, the rational number
\[
 \frac{q^3(\ell+\rho)+Y}{q^3}
\]
has square $\ell(\ell+k)+1\in\Z$.  It is integral, so $q^3\mid Y$.
Consequently
\begin{equation}\label{eq:sc-scale}
 s=Y/q^3>0,\qquad
 \tau=\ell+\rho+s>0,\qquad
 \tau^2-\ell(\ell+k)=1.
\end{equation}
In particular $Y\ge q^3$, $X\ge q$, and the authentic value of $\beta$
gives $q\ge16$.

Put $z_{\rm old}=z+1$.  The transport identity is
\[
 (K_0+w)C+(q-F)-z_{\rm old}(q-1)=1+L_t=1.
\]
It forces $C>0$: if $C\le0$, its left side is at most
\[
 q-F-z_{\rm old}(q-1)
   =1-F-(z_{\rm old}-1)(q-1)\le0.
\]
Hence $\alpha+C=q-F$ and $1\le F\le q-1$.  Set
$G_0=(q-1)(q-F)+\alpha=q\alpha+(q-1)C$.  Positivity gives
\[
 q\le G_0\le q^2-q-1,\qquad
 0<(M_C+qM_F)J<(q-1)(1+2q).
\]
For the upper bound on $G_0$, use $F\ge1$, $C\ge1$, and
$\alpha\le q-F-1$.  Since $J=(q-1)/(B-1)$, the mask bound follows from
the complete tuple's inequalities.  These estimates imply
\begin{equation}\label{eq:sc-bounds}
 3q+1<R<q^4-q^3<q^4-2,
 \qquad E\ge q^4>R+2,\quad a>E>R,
 \quad kY<c<k(Y+1).
\end{equation}

\begin{lemma}[Adaptive modulus]\label{lem:sc-modulus}
After the scale recovery \eqref{eq:sc-scale}, the polynomial $U=Y^2$
satisfies
\[
 U\mid4XY^2,\qquad
 \mu:=\frac{U}{\gcd(U,H)}\ge\frac{Y^2}{9}>R+2.
\]
Neither conclusion requires a value for the main Pell index or a
recovered index quotient.
\end{lemma}
\begin{proof}
The first divisibility holds identically.  Because
$H=4Y(X+1)+3$, we have $\gcd(Y,H)\mid3$.  Prime by prime,
\[
 \min(2v_t(Y),v_t(H))\le2\min(v_t(Y),v_t(H)),
\]
so $\gcd(Y^2,H)\mid9$.  Therefore
\[
 \mu\ge Y^2/9\ge q^6/9>q^4>R+2,
\]
where $q\ge16$ supplies the strict middle inequality.
\end{proof}

The distinction between size and degree is decisive.  Earlier
fourteen-variable quotient absorption used $U=4q^4$.  Here the directly
supplied coordinate $Y$ has degree one; integrality forces $Y\ge q^3$
only on a zero.  Thus $Y^2$ has smaller polynomial degree while meeting
the required arithmetic lower bound.

\subsection{The early bound precedes strong-rank recovery}

The equation $M=0$ makes $g=D_*/U$ a rational square root of the integer
$G$, hence an integer.  It is positive \emph{before} either Pell index
has been identified.  Indeed $r=k-R-1<k$, $a>3$, and $H<5a$, while
\[
 \frac{b_0}{H}>\frac{ac}{H}>\frac c5>\frac{kY}{5}
       >\frac{k}{Y^2}=\frac kU.
\]
The last inequality uses $Y^3>5$.  Hence
\[
 g=b_0+H(\gamma-r/U)>b_0-Hk/U>0,
 \qquad g^2-\Delta c^2=1.
\]
Reducing $D_*=Ug$ modulo $U$ gives
\begin{equation}\label{eq:sc-weakdiv}
 U\mid H(k-R-1),\qquad \mu\mid k-R-1.
\end{equation}
At this stage it is not legitimate to replace $\mu$ by $U$ in the second
divisibility.

Let $A=a+2$ and $P_0=2XY^2+1$.  The first norm is
$\tau^2-XY^2(XY^2+1)k^2=1$.  Its fundamental positive unit has second
coordinate two: coordinate one lies strictly between consecutive
squares, and coordinate two is attained by $2XY^2+1$.  Pell
classification, and then the main norm, give indices $n,p\ge1$ with
\begin{equation}\label{eq:sc-indices}
 \tau=\chi_{P_0}(n),\quad k=2\psi_{P_0}(n),
 \qquad g=\chi_A(p),\quad c=\psi_A(p).
\end{equation}
Since $P_0>A$ and $c>k$, strict growth gives $p>n$.
The recurrence congruence $\psi_{P_0}(n)\equiv n\pmod{P_0-1}$ yields
\begin{equation}\label{eq:sc-firstcong}
 k\equiv2n\pmod{4XY^2}.
\end{equation}
By Lemma~\ref{lem:sc-modulus} and \eqref{eq:sc-weakdiv},
$2n\equiv R+1\pmod\mu$.  Since $0<R+1<\mu$ and $n>0$, we can write
\begin{equation}\label{eq:sc-early}
 2n=R+1+t\mu\quad(t\ge0),\qquad n\ge(R+1)/2\ge25,
 \qquad p\ge26.
\end{equation}
This proves the required early bound without the deleted equation
$k-hE-R=1$.  It also gives
\begin{equation}\label{eq:sc-large}
 c>A\Delta^2,\qquad c>2p,\qquad c>2R.
\end{equation}
For the first inequality, use
$\psi_A(p)\ge(2A-1)^{p-1}\ge A^{25}>A(A^2-1)^2$.
The recurrence gives $c>2p$; finally $c\ge\psi_A(2)=2A>2R$.

\subsection{The ordinary rank and the auxiliary index}

We record the exact inherited rank statement, because a weaker
square-root observation would not suffice.

\begin{lemma}[Ordinary strong rank]\label{lem:sc-rank}
Let $A\ge2$, $\Delta=A^2-1$, and
$c=\psi_A(p)>A\Delta^2$.  If $i,f>0$ and
$(ic^2)^2=\Delta(f^2-1)$, then for an integer $m>0$,
\[
 f=\chi_A(m),\quad p\mid m,\quad c\mid m,\quad
 I=\Delta\psi_A(m),\quad c^2\mid(f^2-1).
\]
In the situation above, $m\ge c>2p$.
\end{lemma}
\begin{proof}
This is the ordinary strong-rank argument preserved in
\path{upstream/ORDINARY_KERNEL.md}, Section 4.  For completeness, its
essential rank step is recalled.  In the quadratic field, let $F_0$ be
the abscissa of the fundamental positive integer-coordinate Pell unit.
For some $d>0$ and $S=\psi_{F_0}(d)$,
\[
 A=\chi_{F_0}(d),\qquad
 \Delta=(F_0^2-1)S^2,\qquad
 \psi_{F_0}(dp)=Sc.
\]
Pell classification applied to the strong equation, using the squarefree
discriminant, gives $f=\chi_{F_0}(u)$ and
$I=(F_0^2-1)S\psi_{F_0}(u)$ for some $u>0$.
Put $b=\psi_{F_0}(\gcd(u,dp))$.
Strong divisibility and $c^2\mid I$ imply $b\ge c/\Delta$.
Indeed, if $g_1=\gcd(c,\psi_{F_0}(u))$, then
$c/g_1\mid(F_0^2-1)S$, while $b\ge g_1$ because
$c\mid\psi_{F_0}(dp)$.  Thus
$b\ge c/((F_0^2-1)S)\ge c/\Delta$.
If $\gcd(u,dp)$ were a proper divisor of $dp$, then
\[
 2b^2<\psi_{F_0}(2\gcd(u,dp))\le Sc<Ac,
\]
contradicting $c>A\Delta^2$.  Thus $dp\mid u$; write $u=dm$,
so $p\mid m$ and composition gives $I=\Delta\psi_A(m)$.

Write $m=pt$.  From $I=ic^2$ and
$\psi_A(pt)=c\psi_{\chi_A(p)}(t)$, reduction modulo $c$ gives
$c\mid\Delta t$.  Here the factor $\chi_A(p)$ is a unit modulo $c$.
Also
$\psi_A(p)\equiv pA^{p-1}\pmod\Delta$, whence
$g_0:=\gcd(c,\Delta)=\gcd(p,\Delta)$ divides $p$.
Since $c/g_0\mid t$, we obtain $c\mid pt=m$.
Finally
\[
 f^2-1=\Delta\psi_A(m)^2=ic^2\psi_A(m),
\]
which proves the required divisibility.  This proof includes the case
where $A$ is not the fundamental unit for the squarefree discriminant.
\end{proof}

Apply the lemma using $W=0$ and \eqref{eq:sc-large}.  Set
$T_{\rm old}=fT>0$ and $y_{\rm old}=|V+e|$.  Then
\begin{equation}\label{eq:sc-aux}
 V=c(T_{\rm old}f-1)-Rf^2,\qquad
 I^2V^2-(I^2-1)y_{\rm old}^2=1.
\end{equation}
Neither $V$ nor $y_{\rm old}$ vanishes; in particular $y_{\rm old}>0$.
Pell classification with parameter $I$ supplies an index $j>0$ such that
\[
 I|V|=\chi_I(j),\qquad y_{\rm old}=\psi_I(j).
\]
The index $j$ is odd, since even-index abscissas are $\pm1$ modulo $I$.
For $j=2v+1$, the integer polynomial $Q_v$ satisfies
\begin{equation}\label{eq:sc-Q}
 \chi_t(2v+1)=tQ_v(t^2),\qquad
 Q_v(0)=(-1)^v(2v+1),\qquad
 Q_v(1-A^2)=(-1)^v\psi_A(2v+1).
\end{equation}
One may verify these identities from
$Q_0(X)=1$, $Q_1(X)=4X-3$, and
$Q_{v+1}(X)=(4X-2)Q_v(X)-Q_{v-1}(X)$.

The explicit expression for $V$ gives
\[
 V\equiv-c\pmod f,\qquad V\equiv-R\pmod c;
\]
the second uses $c^2\mid f^2-1$, which has now been proved by rank
recovery.  Since $I^2\equiv1-A^2\pmod f$, identities
\eqref{eq:sc-Q} imply $\psi_A(j)\equiv\pm c\pmod f$.
Squaring and using the doubling identity gives
\[
 \chi_A(2j)\equiv\chi_A(2p)\pmod{\chi_A(m)}.
\]
The signed Pell step-down therefore yields $j\equiv\pm p\pmod m$.
To see the required uniqueness, reduce the abscissa index modulo $2m$
and reflect it into $[0,m]$ using
\[
 \chi_A(t+2m)\equiv-\chi_A(t),\qquad
 \chi_A(2m-t)\equiv-\chi_A(t)\pmod{\chi_A(m)}.
\]
Because $0<2p<m$, a negative representative cannot have the displayed
positive residue: for $r,2p\le m-1$,
$\chi_A(r)+\chi_A(2p)\le2\chi_A(m-1)<\chi_A(m)$.
The representative $r=m$ has residue zero, and strict growth then forces
the positive representative to be $2p$.  Thus
$2j\equiv\pm2p\pmod{2m}$, as asserted.

Since $c\mid m$, this gives $j\equiv\pm p\pmod c$.
On the other hand $c\mid I$ and the constant term in
\eqref{eq:sc-Q} give $j\equiv\pm R\pmod c$.
Therefore $R\equiv\pm p\pmod c$.  The strict bounds
$0<p,R<c/2$ in \eqref{eq:sc-large} force
\begin{equation}\label{eq:sc-pR}
 p=R.
\end{equation}
This argument uses neither the missing index quotient nor a restored
main-root congruence.

\subsection{Restore the full modulus and finish soundness}

We now have $n<p=R$.  Hence
$|2n-R-1|<R<\mu$.  Its divisibility by $\mu$ in
\eqref{eq:sc-early} forces $2n=R+1$.
Return to the stronger first-Pell congruence
\eqref{eq:sc-firstcong}: it shows that
\[
 r=k-R-1=k-2n
\]
is divisible by the \emph{full} $4XY^2$, and therefore by both $E$ and
$U$.  Moreover $r>0$, since $n\ge2$ and $\psi_{P_0}(n)>n$.
Thus
\begin{equation}\label{eq:sc-recover}
 h=r/E>0,\qquad \gamma_{\rm old}=\gamma-r/U\in\Z.
\end{equation}
The equality $D_*=Ug$ now restores
$g=b_0+H\gamma_{\rm old}$, with the full modulus $H$.
Since $n\ge2$, $k\ge4P_0>X$ and $c>k>X$.
Also $\Delta>(a+1)^2$ implies $g>(a+1)c$, so
\[
 g-b_0>c-X>0,\qquad \gamma_{\rm old}>0.
\]
All fifteen positive coordinates in \eqref{eq:sc-source} have been
reconstructed: use \eqref{eq:sc-scale}, $z_{\rm old}=z+1$,
$T_{\rm old}=fT$, $y_{\rm old}=|V+e|$, and
\eqref{eq:sc-recover}.  Their six source equations hold exactly.
The inherited compiler interface therefore gives a finite ZFC refutation.

The order of the argument is essential.  Initially only
$U\mid Hr$ and $\mu\mid r$ were available.  The bound on $\mu$ first
made the rank theorem applicable; rank recovery gave $p=R$; only then
did the first norm restore $U\mid r$ and the full main-root congruence.

\subsection{Canonical completeness and positive replacement witnesses}

Conversely suppose a finite ZFC refutation exists.  Choose canonical
normalized compiler witnesses as in
\eqref{eq:sc-normalized}--\eqref{eq:sc-canonical}, and replace the strong
coordinate by $i=\Delta i_{\rm norm}$.  Retain the outer coordinates.
The identity $R=2n-1$ and \eqref{eq:sc-firstcong} show that
$r=k-R-1>0$ is divisible by $4XY^2$, hence by $U=Y^2$.
Set $\gamma=\gamma_{\rm old}+r/U>0$.  Then $D_*=Ug$ and $M=0$.

The replacement for the first root also has a positive gap.  Since
$\ell=XY^2k>k$, the equation
$\tau^2=\ell(\ell+k)+1$ gives
\[
 \tau<\ell+k,
 \qquad \tau-\ell=\frac{\ell k+1}{\tau+\ell}>\frac k3.
\]
Canonical $n\ge2$ gives $k>XY$, and therefore
\[
 k/3>XY/3>Y/q^3=s.
\]
Thus $\rho=\tau-\ell-s$ is a positive integer.  Supplying $Y=sq^3$
directly makes $A_0=0$.

The transport shift is positive too.  Since $X=2^R$ and $R>3q$,
$w=2^R/q>q>F$.  As $C>0$, the old transport equation gives
\[
 (z_{\rm old}-1)(q-1)=(K_0+w)C-F>0.
\]
Hence $z=z_{\rm old}-1>0$ and $L_t=0$.

It remains to construct the auxiliary quotient and its positive gap;
an arbitrary old quotient need not be divisible by $f$.
With $I=\Delta i_{\rm norm}c^2$, choose fresh auxiliary coordinates by
\[
 V=\frac{\chi_I(R)}I,
 \qquad y_{\rm old}=\psi_I(R).
\]
The quotient is a positive integer because $R$ is odd.
Write $R=2v+1$; the canonical phase $R\equiv3\pmod4$ makes $v$ odd.
Equation \eqref{eq:sc-normalized} gives
$I^2=\Delta(f^2-1)$ and hence
$I^2\equiv1-A^2\pmod{f^2}$.  Applying the polynomial identities
\eqref{eq:sc-Q} modulo the \emph{square} modulus $f^2$ and modulo $c$
yields
\[
 V\equiv-c\pmod{f^2},\qquad V\equiv-R\pmod c.
\]
Also $\gcd(c,f)=1$ and $f^2\equiv1\pmod c$.
It follows that
\begin{equation}\label{eq:sc-Tconstruct}
 T=\frac{V+c+Rf^2}{cf^2}
\end{equation}
is a strictly positive integer, and its defining identity is exactly
$V=f^2(cT-R)-c$.

The Pell identity gives
$I^2V^2-(I^2-1)y_{\rm old}^2=1$.
Since $R\ge3$ and $I>1$,
$V\ge\chi_I(3)/I=4I^2-3>1$.  Thus
\[
 (I^2-1)(y_{\rm old}^2-V^2)=V^2-1>0,
\]
so $e=y_{\rm old}-V>0$.
Now $y=V+e=y_{\rm old}$, $W=0$, and $N_a=1$.
All fourteen supplied coordinates are positive; Lemma~\ref{lem:sc-minimum}
gives $P=0$.  This proves the converse and completes the semantic part of
Theorem~\ref{thm:sc-main}.

The use of the phase $R\equiv3\pmod4$ in this direction is justified
by the authentic native compiler.  In particular its low-mask condition
$(C-1)\mathbin{\mathrm{AND}}(M_CJ+1)=0$, together with
$q\equiv0\pmod4$, $J\equiv1\pmod4$, and $M_C\equiv2\pmod4$,
gives $C\equiv1\pmod4$.  Substituting in the packing expression for $R$
then gives $R\equiv C+2\equiv3\pmod4$.  This bit condition belongs to
the inherited proof, not the polynomial definition.

\subsection{Exact degree, evaluator cost, and comparison}

Assign degree one to each supplied coordinate and degree zero to fixed
integers.  Put
\[
 q_0=\beta J,\quad X_0=wq_0,\quad k_0=\eta+\zeta,
 \quad E_0=X_0Y,\quad c_0=k_0Y,\quad
 R_0=q_0^3(q_0-F).
\]
These are the relevant leading homogeneous forms.  The exact residual
degrees are
\[
 \deg A_0=12,\qquad \deg M=14,\qquad \deg W=10,
 \qquad \deg N_a=15,\qquad \deg L_t=2.
\]
The cancellations behind the upper bounds follow from exact polynomial
identities:
\begin{align*}
 A_0={}&q^6(2\ell\rho+\rho^2-\ell k-1)
           +2q^3Y(\ell+\rho)+Y^2,\\
 D_0={}&b_0+H\gamma,\\
 D_0^2-G={}&2ac(X+H\gamma)+(X+H\gamma)^2-Hc^2-1,\\
 M={}&U^2(D_0^2-G)-2UHrD_0+H^2r^2,\\
 N_a={}&V^2-(A_{\rm aux}-1)e(2V+e).
\end{align*}
For $M$, the first summand has degree at most $13$, and the other two
have degree $14$.  Their joint leading form is
\[
 \operatorname{in}_{14}M
       =8X_0^2Y^2R_0(k_0Y^3+2R_0),
\]
which is nonzero.  The leaders of $V$ and $A_{\rm aux}$ are
$-f^2R_0$ and $E_0^2f^2$, so
\[
 \operatorname{in}_{15}N_a=2E_0^2f^4eR_0.
\]
Only $N_a^2$ reaches degree $30$.  Therefore
\begin{equation}\label{eq:sc-leader}
 \operatorname{in}_{30}P
  =4\beta^{10}J^{10}w^4Y^4f^8e^2(\beta J-F)^2.
\end{equation}
For example the coefficient of $J^{12}w^4Y^4f^8e^2$ is
$4\beta^{12}\ne0$.  This proves exact degree $30$ with the actual
fixed coefficients, rather than only a syntactic degree bound.

The supplied evaluator uses $44+23+16=83$ binary operations.
Squaring and multiplication by a numeral are charged; reuse and fixed
leaves are free.  All fourteen witness ports, all four fixed ports, and
all $83$ gates are live.  This is an arithmetic-circuit count; it is not a
claim about bit complexity or the size of the fully expanded polynomial.

For orientation, the relevant pinned comparisons are
\begin{center}
\begin{tabular}{lrrr}
\hline
Construction & Positive variables & Degree & Operations\\
\hline
Package 55, fourteen-variable gap form & $14$ & $52$ & $79$\\
Package 58, free-ordinate form & $15$ & $30$ & $79$\\
This construction & $14$ & $30$ & $83$\\
Package 51, shortest evaluator & $15$ & $157$ & $66$\\
\hline
\end{tabular}
\end{center}
Against the previous fourteen-variable point, the degree decreases by
$22/52=42.3077\ldots\%$ for four additional gates.  Against the
fifteen-variable degree-$30$ point, it deletes one positive unknown for
four gates.  These are distinct degree, variable-count, and evaluator-cost
tradeoffs; the result makes no global minimality claim.

\subsection{Provenance and independently reproducible checks}

The comparison is pinned to ProveIt commit
\begin{center}
\path{beb0192cb8dc36f78c369c080f2adfbc1887a486}.
\end{center}
\path{SOURCE_PINS.json} records forty exact copied upstream files, their
repository paths and revisions, Git blob identities, and SHA256 hashes.
The earlier source revision required by the compact coefficient generator
is recorded separately in \path{native/SOURCE_PINS.json}.  The original
notices and license files accompany the native material.

The inherited ingredients are the compact native compiler (package 44),
ordinary strong-rank and signed auxiliary-index arguments, index-quotient
absorption (package 50), the auxiliary gap minimum (package 55), the
squared-modulus quotient (package 51), and integral-root scale recovery
(package 58).  The new ingredient is the choice $U=Y^2$, its early gcd and
size lemma, and its integration with those changes.  Its proof in
Lemma~\ref{lem:sc-modulus} does not assume a decoded power-of-two value
for $X$ or a previously recovered index quotient.

\path{verify.py} is independent of the circuit builder.  It checks source
and coefficient hashes, circuit wiring and liveness, exact local sparse
polynomial identities, independent scalar transcriptions, and entire
univariate polynomial evaluations of the circuit modulo three primes.
It also verifies the exact formal leading forms used in
\eqref{eq:sc-leader}.  Both ordinary and optimized Python executions pass.
The results are recorded in \path{CHECKS.json}.  The complete native
coefficient generator and independent auditor were also rerun in an
isolated copy; all four raw payload hashes agree with the delivered
integers, as recorded in \path{NATIVE_REBUILD.json}.

These finite checks authenticate the delivered implementation and the
explicit finite certificate data.  The universal existence equivalence
rests on the mathematical proof above and the stated compiler interface.
No finite check here decides whether a positive zero or a ZFC refutation
exists.
